\documentclass[conference]{IEEEtran}
\IEEEoverridecommandlockouts

\usepackage[utf8]{inputenc}
\usepackage[T1]{fontenc}
\usepackage[portuguese]{babel}
\usepackage{cite}
\usepackage{amsmath,amssymb,amsfonts}
\usepackage{graphicx}
\usepackage{textcomp}
\usepackage{xcolor}
\usepackage{booktabs}
\usepackage{hyperref}

\def\BibTeX{{\rm B\kern-.05em{\sc i\kern-.025em b}\kern-.08em
    T\kern-.1667em\lower.7ex\hbox{E}\kern-.125emX}}

\begin{document}

\title{Image Captioning com Modelos Multimodais:\\
CLIP + Projetor + SmolLM2}

\author{
\IEEEauthorblockN{Bruno Rodrigues}
\IEEEauthorblockA{\textit{Licenciatura em Engenharia Informática e de Computadores} \\
\textit{ISEL -- Instituto Superior de Engenharia de Lisboa}\\
Lisboa, Portugal \\
a52323@alunos.isel.pt}
}

\maketitle

\begin{abstract}
Neste trabalho implementamos uma pipeline multimodal de \textit{image captioning} que integra um encoder visual pré-treinado com um Large Language Model (LLM) através de um projetor treinável. Usamos o CLIP ViT-B/32 como encoder visual congelado, o SmolLM2-135M-Instruct como LLM congelado, e comparamos dois projetores: uma camada linear e um MLP de duas camadas com ativação GELU. O único módulo treinado é o projetor, com apenas 295K--1,1M parâmetros. A pipeline funciona exclusivamente em CPU, usando o dataset Flickr8k com um subconjunto de 2.000 imagens de treino. Os resultados mostram que o MLP converge mais rapidamente e atinge menor \textit{loss} de validação (3,43 vs 4,50), com ROUGE-L de 0,10 vs 0,05 do projetor linear. O BLEU-4 é nulo em ambos os casos, refletindo as limitações de treinar apenas o projetor com dados reduzidos e um LLM com \textit{priors} conversacionais. O trabalho demonstra a viabilidade da arquitetura e discute em detalhe as causas das limitações observadas.
\end{abstract}

\begin{IEEEkeywords}
image captioning, modelos multimodais, CLIP, projetor, language model, SmolLM2, Flickr8k
\end{IEEEkeywords}


% =============================================================================
\section{Introdução}
% =============================================================================

A geração automática de legendas para imagens (\textit{image captioning}) é uma das tarefas mais representativas na intersecção entre visão computacional e processamento de linguagem natural. O objetivo é produzir uma descrição textual coerente e relevante a partir do conteúdo visual de uma imagem --- uma tarefa que exige compreensão simultânea de ambas as modalidades.

Nos últimos anos, a abordagem dominante evoluiu de modelos CNN+RNN~\cite{b4} para arquiteturas que integram encoders visuais pré-treinados com Large Language Models (LLMs), como exemplificado pelo LLaVA~\cite{b1} e pelo BLIP-2~\cite{b5}. Estas arquiteturas partilham uma ideia central: em vez de treinar um modelo do zero para ambas as modalidades, reutilizam encoders visuais e LLMs pré-treinados, ligando-os através de um módulo de alinhamento de dimensões --- o \textit{projetor} --- que é o único componente treinável. Esta estratégia reduz drasticamente o número de parâmetros a otimizar e permite aproveitar representações já aprendidas em larga escala.

O principal desafio desta abordagem reside no \textit{alinhamento de espaços de representação}: o encoder visual e o LLM foram treinados de forma independente, em dados e objetivos distintos, e os seus espaços de embedding não são diretamente compatíveis. O projetor tem de aprender a traduzir informação visual num formato que o LLM consiga interpretar como contexto para geração de texto.

Neste trabalho, implementamos uma pipeline multimodal leve seguindo este paradigma, com a restrição de funcionar exclusivamente em CPU. A modalidade escolhida é a imagem, e a tarefa é \textit{image captioning} no dataset Flickr8k. A arquitetura consiste em:
\begin{itemize}
    \item \textbf{Encoder visual:} CLIP ViT-B/32~\cite{b3}, congelado, que produz embeddings visuais de 512 dimensões;
    \item \textbf{Projetor:} uma camada linear ou MLP que alinha o espaço visual (512 dims) com o espaço de embeddings do LLM (576 dims);
    \item \textbf{LLM:} SmolLM2-135M-Instruct, congelado, que recebe o embedding visual como prefixo e gera a legenda autoregressivamente.
\end{itemize}

O objetivo principal não é atingir o estado da arte, mas demonstrar que a arquitetura funciona e compreender o papel de cada componente --- em particular, comparar a expressividade de um projetor linear com um MLP de duas camadas, e identificar as limitações desta abordagem quando aplicada com recursos computacionais reduzidos.


% =============================================================================
\section{Trabalhos Relacionados}
% =============================================================================

\subsection{Image Captioning Clássico}

\textbf{Show and Tell}~\cite{b4} (Vinyals et al., 2015) foi um dos primeiros trabalhos a enquadrar \textit{image captioning} como um problema de tradução sequência-a-sequência: uma CNN (GoogLeNet) extrai \textit{features} visuais que inicializam um LSTM gerador de texto. O modelo é treinado de ponta a ponta com cross-entropy sobre as legendas. Esta abordagem estabeleceu o paradigma encoder-decoder e demonstrou que transferência de aprendizagem visual para geração de linguagem era viável, embora o módulo linguístico fosse treinado de raiz e não beneficiasse de LLMs pré-treinados.

\subsection{Encoders Visuais Pré-treinados com Treino Contrastivo}

\textbf{CLIP}~\cite{b3} (Radford et al., 2021) introduziu o treino contrastivo em 400 milhões de pares imagem-texto recolhidos da internet (\textit{WebImageText}), produzindo um encoder visual com representações ricas e transferíveis. O objetivo de treino é maximizar a similaridade coseno entre embeddings de imagens e legendas correspondentes, e minimizá-la para pares não correspondentes. O resultado é um espaço de representação conjunto onde imagens e textos semanticamente relacionados estão próximos. É o encoder que usamos neste trabalho: as suas \textit{features} são suficientemente informativas para \textit{image captioning} sem qualquer fine-tuning adicional.

\subsection{Integração de Visão com LLMs via Projetor}

\textbf{LLaVA}~\cite{b1} (Liu et al., 2023) é a referência mais direta para a nossa arquitetura. Usa CLIP ViT-L/14 como encoder visual e um projetor linear para mapear os embeddings visuais para o espaço de embeddings de um LLM (LLaMA/Vicuna). O treino é feito em duas fases: na primeira, treina apenas o projetor em pares imagem-texto de descrição simples; na segunda, faz \textit{instruction tuning} de todo o modelo em dados de conversação multimodal. A nossa implementação é conceitualmente idêntica à primeira fase do LLaVA, com a diferença de que o LLM permanece congelado durante todo o processo.

\textbf{BLIP-2}~\cite{b5} (Li et al., 2023) propõe uma alternativa ao projetor simples: o Q-Former (\textit{Querying Transformer}), um transformer ligeiro com um conjunto fixo de \textit{query tokens} aprendíveis que interrogam as \textit{features} visuais e produzem um número fixo de tokens visuais para o LLM. O Q-Former atua como gargalo de informação, comprimindo a representação visual e tornando-a mais compatível com o LLM. Esta abordagem é mais sofisticada que um projetor linear mas evita o custo de treinar o LLM completo.

\textbf{LLaMA-VID}~\cite{b2} (Li et al., 2023) estende a ideia do projetor para vídeo, enfrentando o desafio de representar sequências de imagens de forma eficiente. A solução proposta representa cada \textit{frame} com apenas 2 tokens no LLM --- um token de contexto e um token de conteúdo ---, permitindo processar vídeos longos sem explodir o comprimento do contexto. Mostra que a compressão extrema da informação visual é possível mesmo com projetores simples, desde que a informação seja adequadamente selecionada.


% =============================================================================
\section{Dataset}
% =============================================================================

\subsection{Descrição}

Usamos o dataset \textbf{Flickr8k}, disponível no HuggingFace (\texttt{jxie/flickr8k}). Contém 8.000 imagens de situações quotidianas --- crianças a brincar, cães a correr, pessoas em atividades ao ar livre --- cada uma anotada por 5 anotadores humanos independentes, totalizando 40.000 descrições em inglês. As legendas são tipicamente frases curtas de 10 a 15 palavras, descrevendo a ação principal e os sujeitos presentes na imagem. A diversidade de ângulos de descrição entre os 5 anotadores é aproveitada na avaliação: ao usar todas as referências, capturamos melhor a variabilidade natural da linguagem humana.

\begin{table}[htbp]
\caption{Divisões do Dataset Flickr8k}
\begin{center}
\begin{tabular}{lcc}
\toprule
\textbf{Split} & \textbf{Total disponível} & \textbf{Usado (subset CPU)} \\
\midrule
Treino      & 6.000 & 2.000 \\
Validação   & 1.000 & 400   \\
Teste       & 1.000 & 400   \\
\bottomrule
\end{tabular}
\label{tab:dataset}
\end{center}
\end{table}

\subsection{Pré-processamento}

O dataset é carregado diretamente do HuggingFace com a divisão predefinida. As 5 legendas por imagem estão disponíveis nas colunas \texttt{caption\_0} a \texttt{caption\_4}. Para treino, usamos apenas \texttt{caption\_0} por imagem, seguindo a prática comum de usar uma única referência por exemplo de treino. Na avaliação BLEU e ROUGE, todas as 5 legendas são usadas como referências --- prática padrão na literatura para Flickr8k, que permite avaliar cobertura face à variabilidade humana.

O pré-processamento das imagens é feito pelo \texttt{CLIPProcessor}: redimensionamento para $224\times224$ pixels e normalização com as médias e desvios padrão específicos do CLIP (média $[0.481, 0.457, 0.408]$, desvio $[0.269, 0.261, 0.276]$). As \textit{features} CLIP são pré-calculadas uma única vez e guardadas em cache num ficheiro \texttt{.pkl} para eliminar o custo de inferência do encoder em cada epoch de treino.

O tokenizador do SmolLM2 é usado para converter as legendas em sequências de tokens de comprimento máximo 64, com \textit{padding} à direita e truncagem quando necessário. As sequências incluem os tokens especiais \texttt{<bos>} e \texttt{<eos>} para delimitar o início e o fim da legenda.


% =============================================================================
\section{Implementação}
% =============================================================================

\subsection{Encoder Visual}

Usamos o CLIP ViT-B/32 pré-treinado (\texttt{openai/clip-vit-base-patch32}), com 151,3M parâmetros totalmente congelados. A arquitetura ViT-B/32 divide a imagem em patches de $32\times32$ pixels, resultando em $(224/32)^2 = 49$ patches, cada um projetado para um vetor de 768 dimensões. O \textit{transformer} visual processa estes patches com 12 camadas de \textit{self-attention} e produz um vetor \texttt{[CLS]} que representa globalmente a imagem.

As \textit{features} são extraídas através do caminho \texttt{vision\_model} $\rightarrow$ \texttt{pooler\_output} $\rightarrow$ \texttt{visual\_projection}, onde a projeção visual do CLIP reduz as 768 dimensões internas para 512, alinhando o espaço de representação com o encoder de texto do CLIP. O resultado é um vetor $\mathbf{v} \in \mathbb{R}^{512}$ por imagem.

\subsection{Projetor}

O projetor é o único módulo treinável. Mapeia o espaço do CLIP ($d_{clip} = 512$) para o espaço de embeddings do SmolLM2 ($d_{llm} = 576$). A incompatibilidade de dimensões ($512 \neq 576$) impede a utilização direta das \textit{features} CLIP como input do LLM, tornando o projetor necessário mesmo numa versão mínima. Testamos duas versões com complexidades diferentes:

\textbf{Projetor Linear:}
\begin{equation}
    f_L(\mathbf{v}) = \mathbf{W}\mathbf{v} + \mathbf{b}, \quad \mathbf{W} \in \mathbb{R}^{576 \times 512}
\end{equation}
com $576 \times 512 + 576 = 295.488$ parâmetros treináveis. É a solução mais simples possível --- uma transformação linear que roda e escala o espaço visual.

\textbf{Projetor MLP (2 camadas + GELU):}
\begin{equation}
    f_M(\mathbf{v}) = \mathbf{W}_2 \cdot \text{GELU}(\mathbf{W}_1\mathbf{v} + \mathbf{b}_1) + \mathbf{b}_2
\end{equation}
com $\mathbf{W}_1 \in \mathbb{R}^{1024 \times 512}$ e $\mathbf{W}_2 \in \mathbb{R}^{576 \times 1024}$, totalizando 1.115.712 parâmetros. A camada oculta de 1024 unidades com ativação GELU introduz não-linearidade, permitindo ao projetor aprender transformações mais complexas entre os dois espaços. A ativação GELU (\textit{Gaussian Error Linear Unit}) é preferida sobre ReLU em modelos de linguagem por produzir gradientes mais suaves.

\subsection{Language Model}

Usamos o SmolLM2-135M-Instruct (\texttt{HuggingFaceTB/SmolLM2-135M-Instruct}), com 134,5M parâmetros totalmente congelados. O modelo tem \texttt{hidden\_size} = 576, 30 camadas transformer, 9 cabeças de atenção, e foi pré-treinado em dados textuais com fine-tuning de instrução. O modelo é carregado com \texttt{torch\_dtype=torch.float32} para evitar conflitos de tipo com o projetor em ambientes CPU (que não suportam \texttt{bfloat16} de forma nativa).

A integração do sinal visual segue a abordagem do LLaVA: o embedding visual é concatenado como prefixo aos embeddings dos tokens da legenda no espaço de representação do LLM:
\begin{equation}
    \mathbf{E}_{input} = [\mathbf{e}_{vis} \;|\; \mathbf{e}_{t_1}, \mathbf{e}_{t_2}, \ldots, \mathbf{e}_{t_T}] \in \mathbb{R}^{(T+1) \times 576}
\end{equation}
onde $\mathbf{e}_{vis} = f(\mathbf{v}) \in \mathbb{R}^{576}$ é o output do projetor e $\mathbf{e}_{t_i}$ são os embeddings dos tokens de texto obtidos pela camada de embedding do LLM. Esta sequência é passada diretamente ao LLM via \texttt{inputs\_embeds}, contornando a camada de embedding de tokens.

\subsection{Treino}

O objetivo de treino é minimizar a \textit{cross-entropy} causal calculada apenas nas posições dos tokens de texto. O token visual na posição 0 é mascarado com \texttt{label=-100}, indicando ao PyTorch para o ignorar no cálculo do loss. A máscara de atenção cobre a sequência completa (visual + texto), permitindo que o LLM atenda à informação visual ao gerar cada token:
\begin{equation}
    \mathcal{L} = -\frac{1}{|T_{vis}|} \sum_{i=1}^{T} \mathbf{1}[l_i \neq -100] \log p(t_i \mid \mathbf{e}_{vis}, t_1, \ldots, t_{i-1})
\end{equation}

Os hiperparâmetros são apresentados na Tabela~\ref{tab:hparams}. A \textit{learning rate} de $10^{-3}$ é deliberadamente elevada para forçar aprendizagem rápida do projetor em poucas épocas, dado o custo computacional em CPU.

\begin{table}[htbp]
\caption{Hiperparâmetros de Treino}
\begin{center}
\begin{tabular}{ll}
\toprule
\textbf{Parâmetro} & \textbf{Valor} \\
\midrule
Épocas             & 5 \\
Batch size         & 8 \\
Learning rate      & $1 \times 10^{-3}$ \\
Optimizer          & AdamW \\
Max tokens         & 64 \\
dtype              & float32 \\
Caption de treino  & \texttt{caption\_0} \\
Subset de treino   & 2.000 imagens \\
\bottomrule
\end{tabular}
\label{tab:hparams}
\end{center}
\end{table}

\subsection{Inferência}

Na fase de inferência, a geração de legendas é feita de forma autoregressiva com \textit{greedy decoding}. O processo começa com o embedding visual como único contexto e gera tokens um a um até atingir o token \texttt{<eos>} ou o comprimento máximo de 64 tokens:

\begin{enumerate}
    \item Calcular $\mathbf{e}_{vis} = f(\mathbf{v})$ e inicializar a sequência com $[\mathbf{e}_{vis}]$;
    \item Para cada passo $i$: passar a sequência atual ao LLM, obter logits na última posição, selecionar o token com maior probabilidade ($\arg\max$);
    \item Converter o token em embedding e concatenar à sequência;
    \item Parar quando o token gerado for \texttt{<eos>} ou \texttt{<pad>}.
\end{enumerate}

A ausência de \textit{KV-cache} torna a inferência $O(T^2)$ em tempo mas simplifica a implementação. O \textit{greedy decoding} é determinístico e adequado para avaliação comparativa entre projetores.


% =============================================================================
\section{Resultados}
% =============================================================================

\subsection{Curvas de Loss}

A Tabela~\ref{tab:loss} apresenta as curvas de loss de treino e validação para ambos os projetores ao longo das 5 épocas. O MLP converge mais rapidamente logo na primeira época (4,40 vs 5,25 de \textit{train loss}) e continua a melhorar consistentemente, atingindo 3,40 no final do treino contra 4,55 do projetor linear. Em ambos os casos, a \textit{val loss} acompanha a \textit{train loss} sem sinais de \textit{overfitting} significativo, o que é esperado dado o número reduzido de parâmetros treináveis.

\begin{table}[htbp]
\caption{Loss de Treino e Validação por Época}
\begin{center}
\begin{tabular}{ccccc}
\toprule
 & \multicolumn{2}{c}{\textbf{Projetor Linear}} & \multicolumn{2}{c}{\textbf{Projetor MLP}} \\
\cmidrule(lr){2-3} \cmidrule(lr){4-5}
\textbf{Época} & Treino & Val & Treino & Val \\
\midrule
1 & 5,251 & 5,184 & 4,404 & 4,117 \\
2 & 5,209 & 5,282 & 3,879 & 3,652 \\
3 & 5,137 & 4,879 & 3,557 & 3,536 \\
4 & 4,740 & 4,634 & 3,473 & 3,468 \\
5 & 4,553 & 4,500 & 3,404 & 3,434 \\
\bottomrule
\end{tabular}
\label{tab:loss}
\end{center}
\end{table}

\subsection{Métricas de Avaliação}

Usamos duas métricas standard em \textit{image captioning}:

\textbf{BLEU-4} (\textit{Bilingual Evaluation Understudy}) mede a precisão de 4-gramas da legenda gerada face às referências humanas. Um valor de 0 indica que nenhuma subsequência de 4 palavras consecutivas gerada pelo modelo coincide com qualquer referência. É uma métrica severa e mais adequada para modelos que já produzem texto fluente.

\textbf{ROUGE-L} (\textit{Recall-Oriented Understudy for Gisting Evaluation}) mede a maior subsequência comum (\textit{LCS}) entre a legenda gerada e as referências. É mais tolerante a reordenações e menos sensível a fluência do que o BLEU, sendo mais informativa quando o modelo ainda produz texto imperfeito.

\begin{table}[htbp]
\caption{Resultados no Conjunto de Teste (400 imagens, 5 referências por imagem)}
\begin{center}
\begin{tabular}{lcc}
\toprule
\textbf{Projetor} & \textbf{BLEU-4} & \textbf{ROUGE-L} \\
\midrule
Linear & 0,0000 & 0,0521 \\
MLP    & 0,0000 & 0,1023 \\
\bottomrule
\end{tabular}
\label{tab:results}
\end{center}
\end{table}

O BLEU-4 é 0 em ambos os casos, indicando que as legendas geradas não partilham sequências de 4-gramas com as referências humanas. O ROUGE-L mostra que o MLP é aproximadamente $2\times$ melhor que o linear (0,10 vs 0,05), confirmando que a maior expressividade do MLP se traduz em melhores outputs mesmo com métricas de texto.

\subsection{Análise Qualitativa}

A análise das legendas geradas revela dois modos de falha sistemáticos:

\textbf{Degeneração em repetição:} o modelo colapsa para sequências repetitivas, como \textit{``1 2 3 4 5 6 7 8 9 10 ...''} ou repetição de tokens de pontuação. Este comportamento é conhecido na literatura como \textit{repetition loop} e ocorre quando o modelo não consegue escapar de atratores locais no espaço de geração.

\textbf{Texto conversacional:} o LLM, que foi fine-tunado como assistente de instrução, ativa os seus \textit{priors} conversacionais em vez de gerar legendas descritivas. Exemplos observados incluem \textit{``I'm so glad you're doing well! How can I help you today?''} ou \textit{``Sure! Here's a summary of what you asked for.''} --- respostas que fazem sentido num contexto de chat mas não como descrições de imagens.

Nenhum dos projetores gerou legendas que descrevessem corretamente o conteúdo visual das imagens do conjunto de teste. O MLP apresenta outputs ligeiramente mais variados, o que explica o ROUGE-L superior, mas sem qualidade semântica.

\subsection{Discussão}

Os resultados refletem três limitações fundamentais desta implementação:

\textbf{1. Dataset de treino insuficiente.} Com apenas 2.000 exemplos de treino, o projetor vê muito poucas amostras para aprender um mapeamento robusto entre o espaço visual e o espaço linguístico. O LLaVA~\cite{b1} usa 595.000 exemplos na fase de alinhamento. A diferença é de $\approx300\times$.

\textbf{2. LLM congelado com \textit{priors} conversacionais.} O SmolLM2-Instruct foi fine-tunado para responder a instruções em formato de chat. Com o LLM congelado, o projetor tem de competir com estes \textit{priors} muito fortes. Sem \textit{instruction tuning} adicional no formato de caption, o LLM interpreta o embedding visual como um prompt de chat e responde em conformidade.

\textbf{3. Representação visual com um único token.} Toda a informação visual está comprimida num único vetor de 576 dimensões. Modelos estado-da-arte usam dezenas a centenas de tokens visuais para preservar informação espacial. Um único token de contexto é insuficiente para descrever cenas complexas.

A diferença de \textit{loss} entre os dois projetores (3,40 vs 4,55) indica que o MLP aprende um mapeamento superior. Com mais dados e \textit{instruction tuning}, esta diferença provavelmente traduziria em ganhos significativos nas métricas.


% =============================================================================
\section{Conclusão}
% =============================================================================

Implementámos uma pipeline multimodal de \textit{image captioning} funcional, com apenas 295K--1,1M parâmetros treináveis, usando CLIP ViT-B/32 como encoder visual congelado e SmolLM2-135M-Instruct como LLM congelado, ligados por um projetor linear ou MLP. O treino corre exclusivamente em CPU com um subconjunto do Flickr8k.

Os resultados quantitativos são fracos (BLEU-4=0, ROUGE-L$\leq$0,10), mas as curvas de \textit{loss} demonstram que o sistema aprende: a \textit{loss} decresce de forma consistente em ambos os projetores, e o MLP converge para valores substancialmente mais baixos. A fraca qualidade das legendas geradas reflete limitações de recursos --- dados insuficientes, LLM congelado com \textit{priors} conversacionais, e representação visual com um único token --- e não uma falha de arquitetura. A abordagem é conceitualmente idêntica à fase 1 do LLaVA~\cite{b1}, que com ordens de grandeza mais dados e parâmetros treináveis atinge o estado da arte.

O principal contributo deste trabalho é a demonstração de que a arquitectura encoder--projetor--LLM é implementável com recursos mínimos, e que o MLP supera o projetor linear mesmo em condições adversas.


% =============================================================================
\section{Trabalho Futuro}
% =============================================================================

Os resultados obtidos apontam claramente para várias direções de melhoria:

\textbf{Fine-tuning parcial do LLM com LoRA.} O problema mais crítico identificado é o conflito entre os \textit{priors} conversacionais do SmolLM2-Instruct e o formato de caption esperado. Uma solução promissora é o \textit{Low-Rank Adaptation} (LoRA), que permite fazer fine-tuning de uma pequena fração dos parâmetros do LLM (tipicamente $<1\%$) com custo computacional reduzido. Isto permitiria ao modelo aprender o formato de legenda sem perder as capacidades linguísticas gerais.

\textbf{Treino com dataset completo e GPU.} O subconjunto de 2.000 imagens é claramente insuficiente. Treinar com as 6.000 imagens completas do Flickr8k, ou idealmente com o Flickr30k (31.783 imagens), num ambiente com GPU, permitiria avaliar se a arquitetura tem potencial real. A diferença de velocidade entre CPU e GPU em operações de matriz é tipicamente $10\text{--}50\times$, tornando o treino com mais épocas e mais dados viável.

\textbf{Múltiplos tokens visuais.} Em vez de comprimir toda a informação visual num único token, seria possível usar os 49 tokens de patch do ViT diretamente (antes do \texttt{pooler\_output}), projetando cada patch para o espaço do LLM. Isto preservaria informação espacial e permitiria ao LLM ``olhar'' para diferentes regiões da imagem ao gerar cada palavra.

\textbf{Beam search na inferência.} O \textit{greedy decoding} tende a colapsar para sequências de alta frequência. O \textit{beam search} com $k=5$ beams exploraria múltiplas hipóteses em paralelo e tenderia a produzir legendas mais fluentes e diversas.

\textbf{Uso do BLIP-2 como projetor.} Substituir o projetor linear/MLP por um Q-Former pré-treinado (como proposto em BLIP-2~\cite{b5}) introduziria um módulo de alinhamento já treinado em larga escala, potencialmente reduzindo a quantidade de dados necessária para fine-tuning.


% =============================================================================
\begin{thebibliography}{00}

\bibitem{b1}
H. Liu, C. Li, Q. Wu, and Y. J. Lee, ``Visual Instruction Tuning,'' in \textit{Advances in Neural Information Processing Systems (NeurIPS)}, 2023. [Online]. Available: \url{https://arxiv.org/abs/2304.08485}

\bibitem{b2}
K. Li, Y. He, Y. Wang, Y. Li, W. Wang, P. Luo, Y. Wang, L. Wang, and Y. Qiao, ``LLaMA-VID: An Image is Worth 2 Tokens in Large Language Models,'' in \textit{Proc. European Conf. Computer Vision (ECCV)}, 2024. [Online]. Available: \url{https://arxiv.org/abs/2311.17043}

\bibitem{b3}
A. Radford, J. W. Kim, C. Hallacy, A. Ramesh, G. Goh, S. Agarwal, G. Sastry, A. Askell, P. Mishkin, J. Clark, G. Krueger, and I. Sutskever, ``Learning Transferable Visual Models From Natural Language Supervision,'' in \textit{Proc. Int. Conf. Machine Learning (ICML)}, 2021. [Online]. Available: \url{https://arxiv.org/abs/2103.00020}

\bibitem{b4}
O. Vinyals, A. Toshev, S. Bengio, and D. Erhan, ``Show and Tell: A Neural Image Caption Generator,'' in \textit{Proc. IEEE Conf. Computer Vision and Pattern Recognition (CVPR)}, 2015. [Online]. Available: \url{https://arxiv.org/abs/1411.4555}

\bibitem{b5}
J. Li, D. Li, S. Savarese, and S. Hoi, ``BLIP-2: Bootstrapping Language-Image Pre-training with Frozen Image Encoders and Large Language Models,'' in \textit{Proc. Int. Conf. Machine Learning (ICML)}, 2023. [Online]. Available: \url{https://arxiv.org/abs/2301.12597}

\end{thebibliography}

\end{document}
